\documentclass[10pt,a4paper]{amsart}
\usepackage[margin=19mm]{geometry}
\usepackage{amsmath,amssymb,mathtools}
\usepackage[scr=rsfs,cal=euler]{mathalfa}
\usepackage{xfrac}
\usepackage[unicode,hidelinks,bookmarks=false]{hyperref}
\newcommand{\ie}{i.e.\ }
\newcommand{\cf}{cf.\ }
\newcommand{\resp}{respectively\ }
\newcommand{\Subsection}{\subsection}
\providecommand{\nobreakdash}{\nobreak\hbox{-}}
\setlength{\emergencystretch}{2em}
\multlinegap0pt
\usepackage{amssymb}
\usepackage{float}
\usepackage[shortlabels]{enumitem}
\newtheorem{lemma}{Lemma}
\newtheorem{theorem}{Theorem}
\newcommand{\abs}[1]{\left|#1\right|}
\newcommand{\dx}{\,\mathrm{d}x}
\newcommand{\norm}[1]{\left\|#1\right\|}
\newcommand{\ptn}{\partial_t}
\title{Maximum principle and local stability for a class of coupled nonlinear thermo--reaction--phase systems}
\author{Gossrin Jean-Marc Bomisso, Ali Ouattara Kouma, Marie Esther Anass\'e}
\date{Source: arXiv:2604.21756v1 (CC BY 4.0)}
\begin{document}
\maketitle

\noindent\emph{Source and licence.} Maximum principle and local stability for a class of coupled nonlinear thermo--reaction--phase systems. Version arXiv:2604.21756v1 (CC BY 4.0), distributed by arXiv under the Creative Commons Attribution 4.0 International licence, http://creativecommons.org/licenses/by/4.0/, as recorded in the arXiv licence metadata for this exact version and shown on the abstract page https://arxiv.org/abs/2604.21756v1. Full licence notice: \texttt{licenses/arxiv-2604.21756-CC-BY-4.0.txt}.\par\smallskip

\noindent\textit{Editorial scope.} Counted unit: the article's theorem thm:stability (source0.tex lines 791--808). The article's complete original proof of it is delivered with it (source0.tex lines 810--1043, one \texttt{proof} environment of 234 lines). No mathematics is written, repaired or re-derived: the statement and the proof are the article's own lines, in the article's own notation, and no retained line is substituted or re-flowed.  Retained uncounted as the source-local context the proof needs: lines 702--715; lines 722--734; lines 770--781.  Nothing was omitted from any retained span, so the package carries no omission record.  No retained line contains a citation, so the wrapper carries no bibliography and no bibliography entry.  No retained line contains a figure environment, an image-inclusion command or drawing code, so no image is delivered and the package declares no asset dependency.  Editorial formatting only, in addition: an AMS \texttt{amsart} preamble that restates the article's own declarations and macros used by the retained lines, namely the environments \texttt{lemma}, \texttt{theorem} on separate counters, together with the standard packages those lines need.  The retained environments print in this document as Lemma 1, Theorem 1 in order of appearance, whereas the article prints the same items with the numbering of its own full document.  The complete native source of the article as distributed in the arXiv e-print for this exact version is bundled unchanged as \texttt{sources/199047/source0.tex}.
	\begin{equation}
		\label{eq:uvw}
		\left\{
		\begin{aligned}
			\rho c_p \ptn u - \nabla\cdot(k(\phi)\nabla u)
			&= \alpha L_c \ptn v + \alpha L_\phi \ptn w, \\[4pt]
			\ptn v - d_c \Delta v
			&= \mathcal{R}(\theta,c) - \mathcal{R}(\bar\theta,\bar c), \\[4pt]
			\tau_\phi \ptn w - \varepsilon^2\Delta w
			+ (F'(\phi) - F'(\bar\phi))
			&= \lambda(v - w),
		\end{aligned}
		\right.
	\end{equation}

	\begin{equation}
		\label{eq:energy}
		\begin{aligned}
			\mathcal{E}(t) & = \;
			\frac{\rho c_p}{2}\norm{u(t)}_{L^2(\Omega)}^2
			+ \frac12\norm{v(t)}_{L^2(\Omega)}^2
			+ \frac{\tau_\phi}{2}\norm{w(t)}_{L^2(\Omega)}^2
			+ \frac{\varepsilon^2}{2}\norm{\nabla w(t)}_{L^2(\Omega)}^2 \\
			&\quad + \int_\Omega G(w(t))\dx
			- \alpha L_c \int_\Omega u\,v\dx
			- \alpha L_\phi \int_\Omega u\,w\dx,
		\end{aligned}
	\end{equation}

	\begin{lemma}[\textbf{Lipschitz control of} $\mathcal{R}$]
		\label{lem:lipschitzR}
		Assume that $\mathcal T<T_0$. The function
		$$\mathcal{R}(\theta,c) = -\beta K_d(\theta)\,c + \gamma K_r(\theta)\,(1-c)$$
		is Lipschitz continuous in $(\theta, c)$ on $[\underline\theta_{\mathcal{T}}, \infty) \times [0,1]$: there
		exists $L_R > 0$ such that
		\[
		|\mathcal{R}(\theta,c) - \mathcal{R}(\bar\theta,\bar c)|
		\leq L_R\bigl(|u| + |v|\bigr)
		\quad\text{a.e. in } Q_T.
		\]
	\end{lemma}

	\begin{theorem}[\textbf{Local asymptotic stability}]
		\label{thm:stability}
		Under hypotheses \textup{\textbf{(H7)}}-\textup{\textbf{(H10)}}, there exist $\delta > 0$ and $\kappa > 0$ such
		that if
		\[
		\norm{u_0}_{L^2(\Omega)} + \norm{v_0}_{L^2(\Omega)} + \norm{w_0}_{H^1(\Omega)} \leq \delta,
		\]
		then the corresponding solution satisfies, for all $t \in [0, T]$:
		\[
		\mathcal{E}(t) \leq \mathcal{E}(0)\,e^{-\kappa t}.
		\]
		In particular,
		\[
		\norm{u(t)}_{L^2(\Omega)}^2 + \norm{v(t)}_{L^2(\Omega)}^2 + \norm{w(t)}_{H^1(\Omega)}^2
		\leq C\,e^{-\kappa t},
		\]
		and therefore $(\theta, c, \phi) \to (\bar\theta, \bar c, \bar\phi)$ as $t \to +\infty$.
	\end{theorem}

	\begin{proof}
		The proof proceeds in seven steps. For simplicity, we adopt the notation $\norm{\cdot} = \norm{\cdot}_{L^2(\Omega)}$.
		
		\medskip
		\textbf{Step 1: Thermal estimate.}
		We test the first equation of \eqref{eq:uvw} by $u$:
		\begin{equation}
			\label{eq3}
			\frac{\rho c_p}{2}\frac{d}{dt}\norm{u}^2 + k_*\norm{\nabla u}^2
			\leq \alpha L_c \int_\Omega \ptn v\,u\dx + \alpha L_\phi \int_\Omega \ptn w\,u\dx,
		\end{equation}
		using $k(\phi) \geq k_* > 0$ by \textbf{(H1)}.
		
		\medskip
		\textbf{Step 2: Chemical estimate.}
		We test the second equation of \eqref{eq:uvw} by $v$:
		\[
		\frac12\frac{d}{dt}\norm{v}^2 + d_c\norm{\nabla v}^2
		= \int_\Omega (\mathcal{R}(\theta,c) - \mathcal{R}(\bar\theta,\bar c))\,v\dx.
		\]
		By Lemma~\ref{lem:lipschitzR} and Young's inequality ($\eta > 0$ to be chosen):
		\begin{equation}
			\label{eq:v-est}
			\frac12\frac{d}{dt}\norm{v}^2 + d_c\norm{\nabla v}^2
			\leq \eta\norm{u}^2 + C_\eta\norm{v}^2.
		\end{equation}
		
		\medskip
		\textbf{Step 3: Phase estimate.}
		We test the third equation of \eqref{eq:uvw} by $w$:
		\[
		\frac{\tau_\phi}{2}\frac{d}{dt}\norm{w}^2 + \varepsilon^2\norm{\nabla w}^2
		+ \int_\Omega (F'(\phi)-F'(\bar\phi))\,w\dx
		= \lambda\int_\Omega v\,w\dx - \lambda\norm{w}^2.
		\]
		Under hypothesis $F''(\bar\phi) > 0$, there exists $m_F > 0$ such that for $w$ sufficiently small,
		$\int_\Omega (F'(\phi)-F'(\bar\phi))\,w\dx \geq m_F\norm{w}^2$. By Young's inequality:
		\begin{equation}
			\label{eq:w-est}
			\frac{\tau_\phi}{2}\frac{d}{dt}\norm{w}^2 + \varepsilon^2\norm{\nabla w}^2
			+ \left(m_F + \frac\lambda2\right)\norm{w}^2
			\leq \frac\lambda2\norm{v}^2.
		\end{equation}
		
		\medskip
		\noindent\textbf{Step 4: Time derivative of the cross terms and cancellation.}
		
		We compute the time derivative of the two cross terms of
		$\mathcal{E}(t)$ defined in \eqref{eq:energy}.
		
		\begin{itemize}
			\item[\textbf{(i)}]
			By Leibniz's rule, we have
			\begin{equation}
				\label{eq6}
				\frac{d}{dt}\!\left(-\alpha L_c\int_\Omega u\,v\dx\right)
				= -\alpha L_c\int_\Omega (\ptn u)\,v\dx
				-\alpha L_c\int_\Omega u\,(\ptn v)\dx.
			\end{equation}
			The second term (s.t.) of \eqref{eq6} exactly cancels the contribution
			$+\alpha L_c\int_\Omega (\ptn v)\,u\dx$ appearing on the right-hand side
			of \eqref{eq3}:
			\[
			\underbrace{+\alpha L_c\int_\Omega (\ptn v)\,u\dx}_{\text{Step 1}}
			+
			\underbrace{\left( -\alpha L_c\int_\Omega u\,(\ptn v)\dx \right) }_{\text{s.t. of (4.6)}} 
			= 0.
			\]
			It remains to handle $-\alpha L_c\int_\Omega (\ptn u)\,v\dx$.
			We extract $\ptn u$ from the first equation of \eqref{eq:uvw}:
			\[
			\rho c_p\,\ptn u
			= \nabla\cdot\bigl(k(\phi)\nabla u\bigr)
			+ \alpha L_c\,\ptn v
			+ \alpha L_\phi\,\ptn w,
			\]
			whence, after dividing by $\rho c_p$:
			\begin{align}
				\label{eq7}
				\begin{split}
					-\alpha L_c\int_\Omega (\ptn u)\,v\dx
					&= -\frac{\alpha L_c}{\rho c_p}
					\int_\Omega \nabla\cdot\bigl(k(\phi)\nabla u\bigr)\,v\dx \\
					&\quad
					-\frac{\alpha^2 L_c^2}{\rho c_p}
					\int_\Omega (\ptn v)\,v\dx
					-\frac{\alpha^2 L_c L_\phi}{\rho c_p}
					\int_\Omega (\ptn w)\,v\dx.
				\end{split}
			\end{align}
			For the first term on the right-hand side of \eqref{eq7}, integration by parts under homogeneous
			Neumann conditions gives
			\[
			-\int_\Omega \nabla\cdot\bigl(k(\phi)\nabla u\bigr)\,v\dx
			= \int_\Omega k(\phi)\,\nabla u\cdot\nabla v\dx,
			\]
			and Young's inequality with $\varepsilon_1 > 0$, using
			$k(\phi) \leq k^*$ by \textbf{(H1)}, yields
			\[
			\frac{\alpha L_c}{\rho c_p}
			\abs{\int_\Omega k(\phi)\,\nabla u\cdot\nabla v\dx}
			\leq
			\frac{\alpha L_c\, k^*}{\rho c_p}\norm{\nabla u}\,\norm{\nabla v}
			\leq
			\frac{\varepsilon_1}{2}\norm{\nabla u}^2
			+ \frac{\alpha^2 L_c^2\,(k^*)^2}{2\varepsilon_1\,\rho^2 c_p^2}\norm{\nabla v}^2.
			\]
			For the second term on the right-hand side of \eqref{eq7}, we have
			\[
			-\frac{\alpha^2 L_c^2}{\rho c_p}
			\int_\Omega (\ptn v)\,v\dx
			= -\frac{\alpha^2 L_c^2}{2\rho c_p}\frac{d}{dt}\norm{v}^2,
			\]
			which is absorbed into the term $\dfrac{1}{2}\dfrac{d}{dt}\norm{v}^2$ from \eqref{eq:v-est} since $\alpha < \alpha_0$ implies
			$\dfrac{\alpha^2 L_c^2}{\rho c_p} < \dfrac{1}{2} < 1$.
			For the third term on the right-hand side of \eqref{eq7}, Young's inequality gives
			\[
			\frac{\alpha^2 L_c L_\phi}{\rho c_p}
			\abs{\int_\Omega (\ptn w)\,v\dx}
			\leq
			\frac{\alpha^2 L_c L_\phi}{2\rho c_p}
			\bigl(\norm{\ptn w}^2 + \norm{v}^2\bigr),
			\]
			terms which will be absorbed into the dissipative contributions of
			\eqref{eq:w-est} and into $\mu_4\norm{w}^2$ from Step~5 for small $\alpha$.
			\item[\textbf{(ii)}]
			By a strictly analogous computation, we find
			\begin{equation}
				\label{eq8}
				\frac{d}{dt}\!\left(-\alpha L_\phi\int_\Omega u\,w\dx\right)
				= -\alpha L_\phi\int_\Omega (\ptn u)\,w\dx
				-\alpha L_\phi\int_\Omega u\,(\ptn w)\dx.
			\end{equation}
			The second term on the right-hand side of \eqref{eq8} exactly cancels the contribution
			$+\alpha L_\phi\int_\Omega (\ptn w)\,u\dx$ from \eqref{eq3} in Step 1:
			\[
			\underbrace{+\alpha L_\phi\int_\Omega (\ptn w)\,u\dx}_{\text{Step 1}}
			+
			\underbrace{\left( -\alpha L_\phi\int_\Omega u\,(\ptn w)\dx \right)}_{\text{s.t. of (4.8)}}
			= 0.
			\]
			For the remaining term $-\alpha L_\phi\int_\Omega (\ptn u)\,w\dx$,
			we again substitute $\rho c_p\,\ptn u$ from \eqref{eq:uvw} to obtain
			\begin{align*}
				-\alpha L_\phi\int_\Omega (\ptn u)\,w\dx
				&= -\frac{\alpha L_\phi}{\rho c_p}
				\int_\Omega \nabla\cdot\bigl(k(\phi)\nabla u\bigr)\,w\dx \\
				&\quad
				-\frac{\alpha^2 L_c L_\phi}{\rho c_p}
				\int_\Omega (\ptn v)\,w\dx
				-\frac{\alpha^2 L_\phi^2}{\rho c_p}
				\int_\Omega (\ptn w)\,w\dx.
			\end{align*}
			After integration by parts and Young's inequality with
			$\varepsilon_2 > 0$, we obtain
			\[
			\frac{\alpha L_\phi}{\rho c_p}
			\abs{\int_\Omega k(\phi)\,\nabla u\cdot\nabla w\dx}
			\leq
			\frac{\varepsilon_2}{2}\norm{\nabla u}^2
			+\frac{\alpha^2 L_\phi^2\,(k^*)^2}{2\varepsilon_2\,\rho^2 c_p^2}
			\norm{\nabla w}^2.
			\]
			The term involving $\ptn w$ gives
			\[
			-\frac{\alpha^2 L_\phi^2}{\rho c_p}
			\int_\Omega (\ptn w)\,w\dx
			= -\frac{\alpha^2 L_\phi^2}{2\rho c_p}\frac{d}{dt}\norm{w}^2,
			\]
			which is absorbed into $\dfrac{\tau_\phi}{2}\dfrac{d}{dt}\norm{w}^2$ from
			\eqref{eq:w-est} since $\alpha < \alpha_0$ implies
			$\dfrac{\alpha^2 L_\phi^2}{\rho c_p} < \tau_\phi$.
			The remaining cross term is handled by Young's inequality:
			\[
			\frac{\alpha^2 L_c L_\phi}{\rho c_p}
			\abs{\int_\Omega (\ptn v)\,w\dx}
			\leq
			\frac{\alpha^2 L_c L_\phi}{2\rho c_p}
			\bigl(\norm{\ptn v}^2 + \norm{w}^2\bigr).
			\]
			
			\item[\textbf{(iii)}]
			Gathering the contributions from Steps~(i) and~(ii), we obtain,
			for all $\varepsilon_1, \varepsilon_2 > 0$ and $\alpha$ sufficiently small, that
			\begin{align*}
				\frac{d}{dt}\!\left(
				-\alpha L_c\int_\Omega uv\dx
				-\alpha L_\phi\int_\Omega uw\dx
				\right)
				&\leq
				(\varepsilon_1+\varepsilon_2)\norm{\nabla u}^2
				+ C_{\varepsilon_1}\alpha^2\norm{\nabla v}^2
				+ C_{\varepsilon_2}\alpha^2\norm{\nabla w}^2 \\
				&\quad + C\alpha^2\bigl(\norm{v}^2 + \norm{w}^2\bigr),
			\end{align*}
			where the constants $C_{\varepsilon_i} > 0$ depend only on
			$L_c, L_\phi, k^*, \rho, c_p$.
			We choose $\varepsilon_1 = \varepsilon_2 = \dfrac{k_*}{4}$, so that
			$\varepsilon_1 + \varepsilon_2 = \dfrac{k_*}{2}$, in order to absorb these terms
			into $k_*\norm{\nabla u}^2$ from Step~1. For $\alpha < \alpha_0$, the conditions
			\[
			\frac{\alpha^2 L_c^2\,(k^*)^2}{2\varepsilon_1\,\rho^2 c_p^2} \leq \frac{d_c}{2}
			\qquad\text{and}\qquad
			\frac{\alpha^2 L_\phi^2\,(k^*)^2}{2\varepsilon_2\,\rho^2 c_p^2} \leq \frac{\varepsilon^2}{2}
			\]
			are satisfied, and the contributions of order $\alpha^2$ are absorbed into
			$d_c\norm{\nabla v}^2$ and $\varepsilon^2\norm{\nabla w}^2$ from Steps~2 and~3. This prepares the global energy inequality of Step~5.
		\end{itemize}
		\textbf{Step 5: Global energy inequality.}
		Summing the estimates from Steps~1 to~4, with $\varepsilon_1 = \varepsilon_2 = \dfrac{k_*}{4}$,
		$\eta \leq \dfrac{\mu_1 C_P}{2}$ and $\alpha < \alpha_0$, we obtain:
		\[
		\frac{d}{dt}\mathcal{E}(t)
		+ \mu_1\norm{\nabla u}^2 + \mu_2\norm{\nabla v}^2 + \mu_3\norm{\nabla w}^2 + \mu_4\norm{w}^2
		\leq C\bigl(\norm{u}^2 + \norm{v}^2\bigr),
		\]
		with constants $\mu_i > 0$.
		
		\medskip
		\textbf{Step 6: Closure by Poincaré's inequality.}
		Since $\int_\Omega u\dx = 0$ for all time (conservation of the thermal mean under homogeneous Neumann conditions and $H_{\mathrm{ext}} = 0$), Poincaré's inequality
		then gives $\norm{u}^2 \leq C_P\norm{\nabla u}^2$. For $v$, the reaction term controls
		$\norm{v}^2$ via the stability of $\bar c$. Combining with the coercivity of $\mathcal{E}$:
		\[
		\frac{d}{dt}\mathcal{E}(t) + \kappa\,\mathcal{E}(t) \leq 0
		\]
		for some constant $\kappa > 0$.
		
		\medskip
		\textbf{Step 7: Conclusion by Gronwall.}
		The Gronwall lemma applied to $\psi(t) = \mathcal{E}(t)$ gives
		$\mathcal{E}(t) \leq \mathcal{E}(0)\,e^{-\kappa t}$.
		The equivalence of $\mathcal{E}(t)$ with the norm $\norm{u}_{L^2(\Omega)}^2 + \norm{v}_{L^2(\Omega)}^2 + \norm{w}_{H^1(\Omega)}^2$ completes the proof.
	\end{proof}

\end{document}
